% ============================================================
%  ch01a  §1 引言 + §1.1 几何 Satake
%  原文对应：ICM-arXiv.tex 528--656 行
% ============================================================

\section{从经典朗兰兹到几何朗兰兹对应}

本节回顾几何朗兰兹理论的若干进展：其灵感来自经典理论，另一个重要来源是量子物理。基本想法是范畴化／几何化，即把集合论式的陈述替换为范畴式的相应物：
\begin{equation}\label{categorification}
\text{数}\dashrightarrow\text{向量空间}\dashrightarrow\text{范畴}\dashrightarrow\text{2-范畴}\dashrightarrow\cdots .
\end{equation}
下面用几个重要例子来说明这一过程。

\subsection{几何 Satake}\label{S: geometric Satake}

朗兰兹纲领的起点是 Satake 同构（的 Langlands 解释），其中朗兰兹对偶群神秘地登场。类似地，几何朗兰兹理论的起点是几何 Satake 等价：它是连通约化群的（球）局部 Hecke 叠上反常层范畴，与其对偶群的有限维代数表示范畴之间的张量等价。这是经典 Satake 同构的巨大推广（经由层--函数字典），并且可以说为朗兰兹对偶群（实际上是 $C$-群）何以出现在朗兰兹对应中给出了概念性的解释。

关于对偶群的记号与约定，我们遵循 \cite[Sect.~1.1]{Zhu-Integral-Satake}。设 $G$ 为域 $F$ 上的连通约化群，$(\hat G,\hat B,\hat T,\hat e)$ 为其在 $\bZ$ 上的一个带钉朗兰兹对偶群\zy{「带钉」对应原文 pinned：钉 (pinning) 指对偶群的一组结构 $(\hat B,\hat T,\hat e)$，即 Borel 子群、极大环面，以及与根向量相对应的幂幺元素 $e$；它给出同构的确定化，因而在几何 Satake 中必须一并记录。}。存在有限 Galois 扩张 $\widetilde F/F$，以及由 $\Ga_F$ 在 $G$ 的根数据上的作用诱导的单射 $\xi:\Ga_{\widetilde F/F}\hookrightarrow\Aut(\hat G,\hat B,\hat T,\hat e)$。以 $\LG=\hat G\rtimes\Ga_{\widetilde F/F}$ 记 $G$ 的通常 $L$-群，以 $\CG=\hat G\rtimes(\bG_m\times\Ga_{\widetilde F/F})$ 记 \cite{Zhu-Integral-Satake} 中定义的群，它同构于 Buzzard--Gee 引入的 $G$ 的 $C$-群。记 $d:\CG\to\bG_m\times\Ga_{\widetilde F/F}$ 为投影，其核为 $\hat G$。

设 $F$ 为非 阿基米德 局部域，$\mO$ 为其整数环，$k=\bF_q$ 为其剩余域；即 $F$ 是 $\bQ_p$ 的有限扩张，或同构于 $\bF_q(\!(\varpi)\!)$。设 $\sigma$ 为 $k$ 的\emph{几何} $q$-Frobenius 元素。假设 $G$ 可延拓为 $\mO$ 上的连通约化群（这样的 $G$ 称为非分歧的），并固定这样一个延拓，于是有 $G(\mO)\subset G(F)$，后者通常称为 $G(F)$ 的超特殊子群。以 $G(\mO)$ 的有限指数子群作为单位元的开邻域基，$G(F)$ 便成为局部紧拓扑群。经典球 Hecke 代数是 $G(F)$ 上紧支撑、$G(\mO)$-双不变的 $\bC$-值函数构成的空间，其卷积积为
\begin{equation}\label{E:conv prod alge}
(f*g)(x)=\int_{G(F)}f(y)g(y^{-1}x)dy,
\end{equation}
其中 $dy$ 是使 $G(\mO)$ 的体积为 $1$ 的 Haar 测度。注意若 $f$ 与 $g$ 都取 $\bZ$ 值，则 $f*g$ 也取 $\bZ$ 值。因此 $\bZ$ 值函数构成的子集 $H^{\cl}_{G(\mO)}$ 是一个 $\bZ$-代数\footnote{这里 $(-)^{\cl}$ 表示经典 Hecke 代数，以区别于 \eqref{E: Hecke alg isom} 中提到的导出 Hecke 代数。}。

在对偶一侧，在非分歧假设下 $\Ga_{\widetilde F/F}$ 是由 $\sigma$ 生成的有限循环群。注意 $\hat G$ 通过共轭作用在 $\CG|_{d=(q,\sigma)}$ 上，后者是 $d$ 在 $(q,\sigma)\in\bG_m\times\Ga_{\widetilde F/F}$ 处的纤维。于是经典 Satake 同构给出 $\bZ[q^{-1}]$-代数的典范同构
\begin{equation}\label{Sat isom}
\Sat^{\cl}: \bZ[q^{-1}][\CG|_{d=(q,\sigma)}]^{\hat{G}}\cong H^{\cl}_{G(\mO)}\otimes \bZ[q^{-1}].
\end{equation}

\begin{rmk}
事实上，如 \cite{Zhu-Integral-Satake} 所说明，存在定义在 $\bZ$ 上（不反转 $q$）的 Satake 同构，其中 $C$-群 $\CG$ 被替换为某个仿射幺半群，而 $\CG$ 恰为该幺半群的可逆元群。另一方面，若把基环从 $\bZ[q^{-1}]$ 扩张到 $\bZ[q^{\pm \frac{1}{2}}]$，则可把 \eqref{Sat isom} 改写成同构
\begin{equation}\label{Sat isom classical}
 \bZ[q^{\pm \frac{1}{2}}][\hat{G}\sigma]^{\hat{G}}\cong H^{\cl}_{G(\mO)}\otimes \bZ[q^{\pm \frac{1}{2}}],
\end{equation}
其中 $\hat{G}$ 通过通常的共轭作用在 $\hat{G}\sigma\subset {}^LG$ 上（讨论见 \cite{Zhu-Integral-Satake}）。这是 Satake 同构较为传统的表述，略有非典范之处，但对许多应用已经够用。
\end{rmk}

在几何理论中，人们不再把 $G(F)$ 看作拓扑群、考虑其上 $G(\mO)$-双不变紧支撑函数构成的空间，而是把 $G(F)$ 看作某个代数几何对象，研究其上 $G(\mO)$-双等变层的范畴。本节余下部分允许 $F$ 稍微更一般一些。具体地说，设 $F$ 是关于离散赋值完备的局部域，整数环为 $\mO$，剩余域 $k$ 是特征 $p>0$ 的\emph{完全}域\footnote{若 $\cha F=\cha k$（等特征情形），对 $k$ 的这一假设并不必要。此处加上它，是为了统一处理等特征与混合特征（即 $\cha F\neq \cha k$）两种情形。出于同样的理由，下文即使在等特征情形也在完全代数几何的框架下工作。}。设 $\varpi\in\mO$ 为素元。

我们在完全代数几何的框架下工作。回忆：$k$-代数 $R$ 称为完全的，指 Frobenius 自同态 $R\to R,\ r\mapsto r^p$ 是双射。以 $\Aff_k^\pf$ 记完全 $k$-代数构成的范畴。所谓完全预层（更一般地，完全预叠），指从 $\Aff_k^\pf$ 到集合范畴 $\Sets$ 的函子（更一般地，到空间构成的 $\infty$-范畴 $\Spc$ 的函子）。通常代数几何中的许多构造在此框架下依然可行。例如可以对 $\Aff_k^\pf$ 照常赋予 Zariski、étale 或 fpqc 拓扑，并谈论相应的层与叠；随后可定义完全概形、完全代数空间、完全代数叠等等，即具有某些性质的层（叠）。可以证明，这样定义的完全概形／代数空间范畴与通常意义下的完全概形／代数空间范畴等价。完全代数几何的若干基础见 \cite[Appendix A]{Zhu-Satake-Mixed}、\cite{Bhatt-Scholze-Witt} 以及 \cite[\S A.1]{Xiao-Zhu-Cycle}。

设 $R$ 为完全 $k$-代数，以 $W_\mO(R)$ 记系数在 $\mO$ 中的 $R$ 上 Witt 向量环。若 $\cha F=\cha k$，则 $W_\mO(R)\simeq R[[\varpi]]$；若 $\cha F\neq \cha k$，见 \cite[Sect.~0.5]{Zhu-Satake-Mixed}。当 $R=\overline k$ 时，记 $\mO_{\breve F}=W_\mO(\overline k)$，$\breve F=W_\mO(\overline k)[1/\varpi]$。令 $D_R=\Spec W_\mO(R)$，$D_R^*=\Spec W_\mO(R)[1/\varpi]$；可把它们看作以 $\Spec R$ 为参数的一族（打洞）圆盘。

以 $L^+G$（分别地 $LG$）记 $G$ 的\emph{射流群}（分别地\emph{环路群}）。作为 $\Aff_k^\pf$ 上的预层，
\begin{equation*}\label{E:jet and loop}
L^+G(R)=G(W_\mO(R)),\quad LG(R)=G(W_\mO(R)[1/\varpi]).
\end{equation*}
注意 $L^+G(k)=G(\mO)$，$LG(k)=G(F)$。令
\[
\Hk_G:=L^+G\backslash LG/L^+G
\]
为 $LG$ 对左、右 $L^+G$-作用的 étale 叠商，有时称为 $G$ 的（球）局部 Hecke 叠。作为完全预叠，它把 $R$ 送到三元组 $(\mE_1,\mE_2,\beta)$，其中 $\mE_1,\mE_2$ 是 $D_R$ 上的两个 $G$-挠子，而 $\beta:\mE_1|_{D_R^*}\simeq\mE_2|_{D_R^*}$ 是同构。

设 $\ell\neq p$。高阶范畴论的现代发展允许对任意预叠定义 étale $\bF_\ell$-层构成的 $\infty$-范畴（例如见 \cite{Hemo-Zhu-Unipotent}）。特别地，对 $\La=\bF_\ell,\bZ_\ell,\bQ_\ell$（或这些环的有限扩张），可以定义 $\Hk_G$ 上 $\La$-层构成的 $\infty$-范畴 $\Shv(\Hk_G,\La)$，它是 $G(F)$ 上 $G(\mO)$-双不变函数空间的范畴类比。但在不知道 $\Hk_G$ 的若干几何性质时，关于 $\Shv(\Hk_G,\La)$ 几乎无话可说。关键的几何输入是以下定理。

\begin{thm}\label{affGrass}
令 $\Gr_G:=LG/L^+G$ 为 $LG$ 对（右）$L^+G$-作用的 étale 商，它带有左 $L^+G$-作用。则 $\Gr_G$ 可写成 $L^+G$-稳定子函子的归纳极限 $\varinjlim X_i$，其中每个 $X_i$ 都是完全射影簇，且 $X_i\to X_{i+1}$ 是闭嵌入。
\end{thm}

空间 $\Gr_G$ 通常称为 $G$ 的仿射 Grassmann 簇。等特征情形见 \cite{Beauville-Laszlo, Faltings-Loop}，混合特征情形见 \cite{Zhu-Satake-Mixed, Bhatt-Scholze-Witt}；$\Gr_G$ 中闭子簇的例子见 \cite{Zhu-PCMI, Zhu-CDM}。该定理使得可以在 $\Hk_G$ 上定义可构造层与反常层的范畴，并表述几何 Satake 等价，下面就来讨论这一点。

首先，对 $L^+G$-作用取（左）商给出映射 $\Gr_G\to\Hk_G$。粗略地说，$\Hk_G$ 上的层称为反常的（分别地，可构造的），是指它拉回到 $\Gr_G$ 后来自某个 $X_i$ 上的反常（分别地，可构造）层。于是在 $\Shv(\Hk_G,\La)$ 内部有 $\Hk_G$ 上反常层与可构造层的范畴 $\Perv(\Hk_G,\La)\subset\Shv_{c}(\Hk_G,\La)$。可以把它们看作 $G(F)$ 上 $G(\mO)$-双不变紧支撑函数空间的范畴类比。此外，$\Perv(\Hk_G,\La)$ 是 Abel 范畴，当 $\La$ 是特征零的域时还是半单范畴；它称为 $G$ 的 Satake 范畴。为简单起见，下文假定 $\La$ 是域\footnote{当 $\La=\bZ_\ell$ 时相应表述略复杂一些，因为 $\mA$ 与 $\mB$ 为反常层时 \eqref{E: conv product of sheaves} 的右端未必仍是反常层。}。

卷积积 \eqref{E:conv prod alge} 也有范畴类比，即卷积图
\[
 \Hk_G\times\Hk_G\xleftarrow{\pr}  L^+G\backslash LG\times^{L^+G} LG/L^+G\xrightarrow{m} \Hk_G,
\]
而两个层 $\mA,\mB\in\Shv(\Hk_G,\La)$ 的卷积定义为
\begin{equation}\label{E: conv product of sheaves}
\mA\star\mB:= m_!\pr^*(\mA\boxtimes \mB).
\end{equation}
这一卷积积使 $\Shv(\Hk_G,\La)$ 成为幺半 $\infty$-范畴，其幺半子范畴包括 $\Perv(\Hk_G,\La)\subset \Shv_{c}(\Hk_G,\La)$。

\begin{rmk}
上述把 Satake 范畴构造为幺半范畴的方式，与更传统的做法本质等价：传统做法把 Satake 范畴定义为 $\Gr_G$ 上 $L^+G$-等变反常层的范畴（讲解见 \cite{Zhu-PCMI}）。
\end{rmk}

以 $\Coh(\bB\hat{G}_\La)^{\heart}$ 记 $\La$ 上分类叠 $\bB\hat{G}_\La$ 上凝聚层构成的 Abel 幺半范畴\footnote{在对偶群一侧，我们始终在通常代数几何的框架下工作，因此 $\bB\hat{G}$ 是通常意义下的 Artin 叠。}，它等价于 $\hat{G}$ 在有限维 $\La$-向量空间上代数表示构成的范畴。以下定理通常称为几何 Satake 等价。

\begin{thm}\label{intro:geomSat}
存在幺半 Abel 范畴之间的典范等价
\[
\Sat_G:  \Coh(\bB\hat{G}_\La)^{\heart}\cong \Perv(\Hk_G\otimes \overline k,\La).
\]
\end{thm}

几何 Satake 确实是几何朗兰兹纲领的基石之一，并已在表示论、数学物理以及（算术）代数几何中找到众多应用。当 $F=k(\!(\varpi)\!)$ 时，该定理源自 Lusztig、Ginzburg、Beilinson--Drinfeld 与 Mirković--Vilonen 的工作（参见 \cite{Lusztig-q-analog, Beilinson-Drinfeld-Hitchin, Mirkovic-Vilonen-GeoSat}）。在混合特征情形，它在 \cite{Zhu-Satake-Mixed, Yu-Integral-Satake} 中以等特征情形为输入得到证明；\cite{Fargues-Scholze-Geometrization} 则通过仿照等特征情形的策略给出证明。本节最后给出几条注记。

\begin{rmk}
\label{R: birth of the dual group}
(1) 如前所述，几何 Satake 可看作 $G$ 的朗兰兹对偶群 $\hat G$ 的概念性定义，即 Tannaka 范畴 $\Perv(\Hk_G\otimes \overline k,\La)$ 的 Tannaka 群。此外，如 \cite{Yun-Zhu-Integral-Homology, Zhu-Ramified-Sat} 所说明，$\hat G$ 典范地带有一个钉 $(\hat B, \hat T, \hat e)$。本文余下部分中，所谓 $G$ 的带钉朗兰兹对偶群 $(\hat{G},\hat{B},\hat{T},\hat{e})$，均指由几何 Satake 定义的这个四元组。

(2) 对算术应用而言，需要用对偶群一侧来理解 $\Ga_k$ 在 $\Perv(\Hk_G\otimes\overline k,\La)$ 上的作用。结果表明，这一作用由 $\Ga_k$ 在 $\hat{G}$ 上的某个作用诱导，该作用保持 $(\hat{B},\hat{T})$ 但不保持 $\hat{e}$。见 \cite{Zhu-Ramified-Sat, Zhu-PCMI}，或从 motivic 角度出发的 \cite{Zhu-CDM}。这正是群 $\CG$ 出现的原因。详细讨论见 \cite{Zhu-Ramified-Sat, Zhu-PCMI, Zhu-Integral-Satake}。

(3) 还有导出 Satake 等价 \cite{Bezrukavnikov-Finkelberg-Derived-Satake}，至少当 $\La$ 是特征零的域时，它用对偶群来描述 $\Shv_{c}(\Hk_G\otimes\overline k,\La)$。需要指出，对偶一侧出现的范畴并不是 $\bB\hat{G}_\La$ 上凝聚层的导出范畴。

(4) 事实上，对许多应用来说，重要的是几何 Satake 的族版本。设 $S$ 为（非空）有限集，则在 $D^S$ 上有局部 Hecke 叠 $\Hk_{G,D^S}$，这里 $D^S$ 是圆盘 $D=\Spec \mO$ 在 $S$ 上的自乘积。粗略地说，它分类四元组 $(\{x_s\}_{s\in S}, \mE,\mE',\beta)$，其中 $\{x_s\}_{s\in S}$ 是 $D$ 的一个 $S$-元点组，$\mE$ 与 $\mE'$ 是 $D$ 上两个 $G$-挠子，$\beta$ 是 $\mE$ 与 $\mE'$ 在 $D-\cup_{s} \{x_s\}$ 上的同构。在等特征情形，可把 $D$ 看成 $k$ 上代数曲线 $X$ 在某个 $k$-点处的形式圆盘，而 $\Hk_{G,D^S}$ 就是叠
\[
\Hk_{G, X^S}=(L^+G)_{X^S}\backslash (LG)_{X^S}/(L^+G)_{X^S}
\]
沿 $D^S\to X^S$ 的限制，其中 $(LG)_{X^S}$ 与 $(L^+G)_{X^S}$ 是 $LG$ 与 $L^+G$ 在 $X^S$ 上的族版本（精确定义见 \cite[Sect.~3.1]{Zhu-PCMI}）。在混合特征情形，叠 $\Hk_{G,D^S}$（事实上连 $D^S$ 本身）并不生活在（完全）代数几何的世界中，而生活在 Scholze 所发展的完满解析几何世界中（见 \cite{Scholze-Berkeley, Fargues-Scholze-Geometrization}）。两种情形下都可以考虑 $\Hk_{G,D^S}\otimes\overline k$ 上（ULA）反常层的范畴 $\Perv^{\on{ULA}}(\Hk_{G,D^S}\otimes\overline k,\La)$。此外，对有限集之间的映射 $S\to S'$，沿 $\Hk_{G,D^{S'}}\to\Hk_{G,D^S}$ 的限制给出函子
$\Perv^{\on{ULA}}(\Hk_{G,D^S}\otimes\overline k,\La)\to \Perv^{\on{ULA}}(\Hk_{G,D^{S'}}\otimes\overline k,\La)$\footnote{这个限制函子定义了所谓的融合积，它是几何 Satake 等价中的一个关键概念。``fusion'' 这一术语最初来自共形场论。}。细节参见上述文献。

另一方面，令 $\hat{G}^S$ 为 $\hat{G}$ 在 $\La$ 上的 $S$ 次自乘积。则对 $S\to S'$，沿 $\bB \hat{G}^{S'}\to \bB \hat{G}^S$ 的限制给出函子 $\Coh(\bB \hat{G}_{\La}^S)^\heart\to \Coh(\bB \hat{G}_{\La}^{S'})^\heart$。几何 Satake 的族版本给出函子族
\begin{equation}\label{factorizable Sat}
\Sat_S: \Coh( \bB \hat{G}_{\La}^S)^\heart\to \Perv^{\on{ULA}}(\Hk_{G,D^S}\otimes\overline k,\La),
\end{equation}
它与两侧由有限集之间的映射诱导的限制函子相容（见 \cite{Gaitsgory-deJong, Zhu-PCMI}）。

(5) 对应用而言，重要的是在不同层论框架下、在不同版本的局部 Hecke 叠上都有几何 Satake。除上面提到的之外，我们还要提到 $D$-模版本 \cite{Beilinson-Drinfeld-Hitchin}，以及算术 $D$-模版本 \cite{Xu-Zhu-Bessel}。
\end{rmk}
